\lysh{16191_SOLUTION}{1}
\textbf{$\boldsymbol{\alpha})$}
\quad \text{i.}
\[ \overrightarrow{\mathrm{AM}}=(x-1, \psi-1) \]
\[ \overrightarrow{\mathrm{BM}}=(x-5, \psi-5) \]
\[ \overrightarrow{\mathrm{AM}}^{2}+\overrightarrow{\mathrm{BM}}^{2}=32 \Rightarrow \left(\sqrt{(x-1)^{2}+(\psi-1)^{2}}\right)^{2}+\left(\sqrt{(x-5)^{2}+(\psi-5)^{2}}\right)^{2}=32 \]
\[ x^{2}-2x+1+\psi^{2}-2\psi+1+x^{2}-10x+25+\psi^{2}-10\psi+25=32 \]
$2x^{2}+2\psi^{2}-12x-12\psi+20=0$ διαιρούμε με $2$ και έχουμε
\[ x^{2}+\psi^{2}-6x-6\psi+10=0 \quad (1) \]
\quad \text{ii.} Για να παριστάνει η εξίσωση $x^{2}+\psi^{2}+\mathrm{A}x+\mathrm{B}\psi+\Gamma=0$ κύκλο θα πρέπει $\mathrm{A}^{2}+\mathrm{B}^{2}-4\Gamma>0$. \\
Από την $(1)$ έχουμε $\mathrm{A}=-6, \mathrm{B}=-6, \Gamma=10$.
\[ \mathrm{A}^{2}+\mathrm{B}^{2}-4\Gamma=(-6)^{2}+(-6)^{2}-4\cdot 10 = 36+36-40=32>0 \]
Επομένως πρόκειται περί κύκλου.

\textbf{$\boldsymbol{\beta})$}
\quad \text{i.} Για να εφάπτεται ο κύκλος στην ευθεία, πρέπει η απόσταση του κέντρου $\mathrm{K}$ από την ευθεία να ισούται με την ακτίνα του κύκλου. \\
$d(\mathrm{K},\varepsilon)=\rho$ δηλαδή $\frac{|3\lambda+3-2|}{\sqrt{\lambda^{2}+1}} = 2\sqrt{2} \Rightarrow |3\lambda+1|=2\sqrt{2}\sqrt{\lambda^{2}+1}$
\[ (3\lambda+1)^{2}=8(\lambda^{2}+1) \]
$9\lambda^{2}+6\lambda+1=8\lambda^{2}+8 \Rightarrow \lambda^{2}+6\lambda-7=0$. \\
Υπολογίζουμε τις ρίζες και έχουμε $\lambda=-7$ και $\lambda=1$.
\quad \text{ii.} Ο συντελεστής διεύθυνσης της $\mathrm{AB}$ είναι $\lambda_{\mathrm{AB}}=\frac{\psi_{\mathrm{B}}-\psi_{\mathrm{A}}}{\chi_{\mathrm{B}}-\chi_{\mathrm{A}}}=\frac{5-1}{5-1}=1$. Ένα διάνυσμα παράλληλο στην $\mathrm{AB}$ είναι το $\vec{\delta_{1}}=(1,1)$ ενώ ένα διάνυσμα παράλληλο στην $(\varepsilon)$ είναι το $\vec{\delta_{2}}=(1,-\lambda)$. Η γωνία των δύο ευθειών είναι η γωνία των δύο διανυσμάτων που είναι παράλληλα σε αυτές.
\[ \sigma\upsilon\nu\widehat{\left(\vec{\delta_{1}},\vec{\delta_{2}}\right)}=\sigma\upsilon\nu 45^{\circ} \text{ ή } \frac{\vec{\delta_{1}}\cdot\vec{\delta_{2}}}{|\vec{\delta_{1}}||\vec{\delta_{2}}|}=\frac{\sqrt{2}}{2} \text{ ή } \frac{(1)(1)+(1)(-\lambda)}{\sqrt{1^{2}+1^{2}}\sqrt{1^{2}+(-\lambda)^{2}}}=\frac{\sqrt{2}}{2} \]
\[ \frac{1-\lambda}{\sqrt{2}\sqrt{1+\lambda^{2}}}=\frac{\sqrt{2}}{2} \]
$2(1-\lambda)=\sqrt{2}\sqrt{2}\sqrt{1+\lambda^{2}}$ ή $2(1-\lambda)=2\sqrt{1+\lambda^{2}}$ ή $1-\lambda=\sqrt{1+\lambda^{2}}$
\[ (1-\lambda)^{2}=(\sqrt{1+\lambda^{2}})^{2} \]
$1-2\lambda+\lambda^{2}=1+\lambda^{2}$ ή $\lambda=0$.